\chapter{Conclusion and Future Work}
\label{Conclusion_and_Future_Work}

\section{Summary of Achievements}

This project successfully established a comprehensive framework for a scalable, low-cost Indoor Positioning System based on ESP32 microcontrollers and the ESP-NOW protocol \cite{espressif2023esp32, espressif2023espnow}. The work demonstrates that high-precision localization can be achieved using readily available, inexpensive hardware through a carefully designed phased architectural approach \cite{zafari2019survey}.

The primary achievements of this project include:

\subsection{Phased Architecture Implementation}

The project introduced a three-phase architectural evolution that ensures compatibility and incremental complexity management. Phase 1 (Hands-Free Single Target) establishes the foundational calibration protocol and validates the path-loss model with expected precision of $\pm 5$--$10$ cm. Phase 2 (Multi-Target Sequential) extends the system to handle multiple concurrent targets through time-slot management, maintaining precision within $\pm 10$--$15$ cm. Phase 3 (All-Nodes Mesh Network) provides the ultimate goal of fully distributed autonomous positioning with relative precision of $\pm 20$ cm. This phased approach allows each stage to build upon the previous one while maintaining backward compatibility, ensuring that lessons learned and calibration data from earlier phases remain valid.

\subsection{Calibration Protocol Establishment}

A rigorous calibration protocol was developed to derive environment-specific path-loss constants ($A$ and $n$) from the standard RSSI-distance relationship: $d = 10^{(A - \text{RSSI})/(10n)}$ \cite{rappaport1996wireless, seidel1992path}. The protocol requires collecting 50--100 RSSI samples per anchor at known positions to establish accurate baseline parameters \cite{goldoni2010experimental}. This calibration phase eliminates environmental variables and provides the foundation for all subsequent positioning operations. The protocol emphasizes static node placement during calibration to minimize motion-induced interference and multipath effects \cite{haeberlen2004practical}.

\subsection{Machine Learning Model Baseline}

An MLPRegressor-based machine learning model was developed and evaluated using synthetic data to explore direct RSSI-to-coordinate mapping as an alternative to traditional trilateration \cite{xiao2016indoor, wang2015indoor}. The model architecture consists of two hidden layers with 64 neurons each, using ReLU activation functions and StandardScaler normalization. Initial results on synthetic data demonstrate the feasibility of this approach, with performance metrics (MSE and R² scores) providing a baseline for future improvement with real experimental data \cite{zhang2019deep}. The feasibility analysis identified concrete strategies for enhancement, including increased training data volume, feature engineering techniques, and hyperparameter optimization.

\subsection{Low-Cost Reproducible Approach}

The entire system is designed around commercially available ESP32 development boards, requiring minimal additional hardware (batteries, OLED displays for debugging) \cite{espressif2023esp32}. The testbed operates on a compact 70$\times$50 cm grid, making it suitable for laboratory environments and educational settings. The use of ESP-NOW protocol eliminates the need for complex networking infrastructure, reducing both cost and setup complexity \cite{espressif2023espnow}. This reproducibility ensures that the system can be easily replicated by other researchers or students for further experimentation and validation \cite{sadowski2018rssi}.

\section{Key Takeaways}

The development and analysis of this Indoor Positioning System revealed several critical operational principles that are essential for successful RSSI-based localization:

\subsection{Importance of Static Calibration}

Accurate localization fundamentally depends on establishing a reliable path-loss model before attempting multi-node operations \cite{rappaport1996wireless}. The calibration phase must be conducted with all nodes in static positions to eliminate motion-induced RSSI fluctuations \cite{goldoni2010experimental}. Environmental factors such as multipath fading, temperature variations, and RF interference significantly affect the path-loss constants $A$ and $n$, making environment-specific calibration mandatory rather than optional \cite{lemic2017experimental}. Attempting to skip or shortcut the calibration process inevitably leads to degraded positioning accuracy.

\subsection{Sequential Timing for Collision Avoidance}

In multi-target scenarios, uncoordinated simultaneous broadcasts create signal collisions that corrupt RSSI measurements and make source identification ambiguous \cite{lemic2017experimental}. The sequential time-slot protocol, with carefully managed intervals (e.g., 100 ms spacing), is essential for maintaining data integrity. Each target must broadcast during its designated time slot, with anchors tagging received data with both target ID and timestamp. This approach prevents RF noise floor elevation and ensures that each RSSI measurement can be unambiguously attributed to its source \cite{sadowski2018rssi}.

\subsection{RSSI Consistency Over Sample Count}

Data quality is more important than data quantity for RSSI-based positioning \cite{haeberlen2004practical}. Collecting 50--100 samples per measurement point and applying filtering techniques (moving average or median filtering) produces more reliable results than using thousands of unfiltered instantaneous measurements \cite{goldoni2010experimental}. The inherent variability in RSSI due to multipath effects and environmental noise means that statistical aggregation is necessary to extract meaningful distance information \cite{pahlavan2002indoor}. This principle applies across all three phases of the system architecture.

\subsection{Distributed Mesh Autonomy Benefits}

The All-Nodes Mesh Network architecture represents the only viable path toward fully autonomous positioning without centralized infrastructure \cite{akyildiz2005wireless}. By enabling bidirectional pairwise RSSI measurements between all nodes, the system constructs a comprehensive adjacency map that can be processed using graph optimization techniques \cite{moore2004robust}. This distributed approach eliminates single points of failure, reduces communication bottlenecks, and enables scalability to larger networks. The mesh architecture also supports advanced algorithms such as SLAM-style positioning and multilateration/MDS techniques that leverage the complete network topology \cite{shang2003localization, ji2006sensor}.

\section{Evolution Roadmap}

The three-phase architecture provides a clear roadmap for system evolution, with each stage building upon the achievements of the previous one:

\subsection{Stage 1: Single Target (Static Calibration)}

This initial stage focuses on establishing the fundamental positioning capability with a single mobile target and multiple fixed anchors \cite{bahl2000radar}. The primary objectives are deriving accurate path-loss constants ($A$, $n$) through systematic calibration, validating trilateration algorithms with known ground truth positions, and achieving target precision of $\pm 5$--$10$ cm \cite{cheung1994least}. This stage has been completed in the current project, providing the foundation for subsequent phases.

\subsection{Stage 2: Sequential Multi-Target}

The second stage extends the system to handle multiple concurrent targets through time-slot management. Implementation involves integrating the sequential broadcast protocol, validating collision avoidance mechanisms, and testing with 2--4 simultaneous targets. The target precision for this stage is $\pm 10$--$15$ cm, with the slight degradation compared to Stage 1 attributed to increased system complexity and potential timing synchronization challenges.

\subsection{Stage 3: All-Nodes Mesh Network}

The final stage implements fully distributed positioning where all nodes act as both targets and anchors. This requires bidirectional ESP-NOW communication, construction of complete RSSI adjacency maps, and application of advanced multilateration or graph optimization algorithms. The target relative precision is $\pm 20$ cm, reflecting the increased complexity of distributed computation and the absence of fixed reference anchors with known absolute positions.

\subsection{Model Improvement with Real Data}

As real experimental data becomes available from the physical testbed, the machine learning model can be significantly enhanced \cite{xiao2016indoor}. The synthetic dataset used for baseline evaluation contains only 200 samples, whereas a comprehensive real-world dataset could include thousands of measurements across the entire 70$\times$50 cm grid. Feature engineering techniques such as RSSI ratios between anchor pairs, signal gradient calculations, and temporal consistency metrics can be incorporated to improve model performance \cite{he2016wifi}. Hyperparameter tuning and exploration of ensemble methods (Random Forest, Gradient Boosting) may further enhance accuracy, potentially surpassing traditional trilateration approaches in complex multipath environments \cite{zhang2019deep}.

\subsection{Advanced Algorithms}

Future iterations of the mesh network can incorporate sophisticated positioning algorithms beyond basic trilateration \cite{cheung1994least}. Graph optimization techniques, similar to those used in Simultaneous Localization and Mapping (SLAM), can process the complete adjacency map to refine position estimates iteratively \cite{moore2004robust}. Multidimensional Scaling (MDS) algorithms can transform the pairwise distance matrix into a coherent coordinate system \cite{ji2006sensor, shang2003localization}. These advanced approaches leverage the rich information available in the mesh topology to achieve better accuracy and robustness than single-point trilateration methods \cite{biswas2006semidefinite}.

\section{Future Work}

Several promising directions exist for extending and enhancing the Indoor Positioning System:

\subsection{Real-World Deployment and Testing}

The current system has been designed and validated using synthetic data and theoretical analysis. The next critical step is comprehensive real-world testing on the physical testbed \cite{goldoni2010experimental}. This involves deploying the hardware across the 70$\times$50 cm grid, conducting extensive calibration measurements, collecting training data for the ML model, and validating positioning accuracy against ground truth measurements. Real-world testing will reveal practical challenges such as battery life management, thermal effects on RSSI stability, and the impact of human presence on signal propagation \cite{lemic2017experimental}.

\subsection{Dynamic Environment Adaptation}

Current calibration assumes static environmental conditions, but real-world deployments must contend with dynamic factors such as temperature variations, humidity changes, and moving obstacles \cite{yassin2017recent}. Future work should explore adaptive calibration techniques that can detect environmental changes and adjust path-loss parameters in real-time \cite{lemic2017experimental}. Machine learning approaches could be employed to model the relationship between environmental sensors (temperature, humidity) and RSSI behavior, enabling predictive compensation for environmental effects \cite{zhang2019deep}.

\subsection{Integration with Existing Infrastructure}

The ESP32 platform supports multiple communication protocols including Wi-Fi and Bluetooth in addition to ESP-NOW \cite{espressif2023esp32}. Future implementations could integrate the positioning system with existing wireless infrastructure, using Wi-Fi for data logging and remote monitoring while maintaining ESP-NOW for time-critical positioning operations \cite{jimenez2018indoor}. Integration with cloud services could enable centralized data analysis, remote system management, and multi-site deployment coordination.

\subsection{Scalability to Larger Networks}

While the current testbed uses 4--5 anchor nodes covering a 70$\times$50 cm area, practical applications may require coverage of much larger spaces \cite{liu2007survey}. Scaling to larger networks introduces challenges in time-slot management (more nodes require longer cycle times), communication range limitations, and computational complexity of graph optimization algorithms \cite{akyildiz2005wireless}. Future research should investigate hierarchical network architectures, where local clusters of nodes perform positioning independently and coordinate through gateway nodes to maintain global consistency \cite{moore2004robust}.

\subsection{Machine Learning Model Improvements}

Beyond the baseline MLPRegressor model, several advanced machine learning techniques warrant investigation \cite{zhang2019deep}. Deep neural networks with additional hidden layers and more sophisticated architectures (e.g., residual connections, attention mechanisms) may capture complex nonlinear relationships between RSSI patterns and positions \cite{xiao2016indoor}. Recurrent neural networks (RNNs) or Long Short-Term Memory (LSTM) networks could exploit temporal correlations in sequential RSSI measurements to improve accuracy and reduce sensitivity to instantaneous noise \cite{hoang2019recurrent, chen2018wifi}. Transfer learning approaches could enable models trained in one environment to be adapted to new environments with minimal additional calibration data.

\subsection{Power Optimization}

For battery-powered deployments, power consumption becomes a critical constraint. Future work should investigate duty-cycling strategies where nodes alternate between active and sleep modes, wake-on-radio techniques that minimize idle power consumption, and adaptive transmission power control that balances positioning accuracy against energy efficiency. The ESP32's deep sleep capabilities could be leveraged to extend battery life from hours to days or weeks, making long-term deployments practical.

\subsection{Multi-Floor and 3D Positioning}

The current system focuses on 2D positioning within a single plane. Extending to 3D positioning requires additional anchors placed at different heights and modifications to the trilateration algorithms to solve for three coordinates instead of two \cite{cheung1994least}. Multi-floor positioning in buildings introduces additional challenges due to floor attenuation effects and the need for vertical anchor placement \cite{yassin2017recent}. Future research could explore hybrid approaches combining RSSI-based horizontal positioning with barometric pressure sensors for vertical positioning \cite{jimenez2018indoor}.

\section{Concluding Remarks}

This project demonstrates that accurate, low-cost indoor positioning is achievable using commodity ESP32 hardware and the ESP-NOW protocol \cite{espressif2023esp32, espressif2023espnow}. The phased architectural approach provides a practical roadmap from initial calibration through multi-target operation to fully distributed mesh networks \cite{akyildiz2005wireless}. The combination of traditional trilateration techniques with machine learning approaches offers flexibility to optimize for different deployment scenarios and accuracy requirements \cite{xiao2016indoor, cheung1994least}. The foundation established by this work enables future researchers and developers to build upon these principles and extend the system toward practical real-world applications in asset tracking, robotics navigation, and location-based services \cite{zafari2019survey, yassin2017recent}.
